\section{Introduction}
Surface irrigation spreads water over the soil by gravity: the water advances across the field, soaks in and, after the inflow is cut, recedes or ponds. It irrigates most of the world's irrigated land, because it is cheap. Its performance depends on a small number of things: the intake of the soil, the slope, the roughness, the length, and the inflow and cut-off time. This chapter gives the design of the furrow, border and basin, with a full worked furrow.

\section{Types of surface irrigation}
There are three. In \emph{basins} the water is ponded on level land and the soil drinks it; the design is the volume per basin. In \emph{borders} a strip of land is irrigated by a sheet of water that advances down a gentle slope. In \emph{furrows} the water runs in small channels between crop rows and wets the soil sideways; furrows suit row crops and give the farmer the most control. Basins need flat land and a large flow; borders suit close-growing crops such as wheat and alfalfa; furrows suit maize, cotton and vegetables.

\section{Soil intake}
The cumulative intake per unit length of furrow is $Z=k\tau^{a}+f_0\tau$ (Kostiakov--Lewis), with the opportunity time $\tau$ in minutes. With $k=0.0040$ m$^3$\,m$^{-1}$\,min$^{-a}$, $a=0.55$ and $f_0=0.0001$ m$^3$\,m$^{-1}$\,min$^{-1}$, a required depth of 60 mm over a spacing of 0.75 m means $Z=0.0450$ m$^3$\,m$^{-1}$; solving $k\tau^a+f_0\tau=Z$ gives $\tau_n=62$ min.

\section{Flow in the furrow}
The flow of the furrow follows Manning's equation $Q=\frac{1}{n}AR^{2/3}S_0^{1/2}$. With $n=0.04$, $A=0.012$ m$^2$, $R=0.05$ m and $S_0=0.002$ the carrying capacity is $1.8$ L\,s$^{-1}$; the safe inflow against erosion is Criddle's $q_{max}=0.63/S_0(\%)$ = 3.1 L\,s$^{-1}$ for a slope of 0.2 per cent.

